\documentclass[12pt]{article}
\usepackage{bbm}
\usepackage{fancyhdr}
\usepackage{titling}
\usepackage{enumitem}
\usepackage{amsmath}
\usepackage{dsfont}
\usepackage{amsfonts}
\usepackage{amssymb}
\usepackage{amsthm}
\usepackage{geometry}
\usepackage{tikz}
\usepackage{algorithm}
\usepackage{algpseudocode}
\geometry{verbose,tmargin=1in,bmargin=1in,lmargin=1in,rmargin=1in}
\renewcommand{\vec}[1]{\mathbf{#1}}
\newcommand{\dec}{\text{dec}}
\newcommand{\divi}{\mathbin{\textnormal{div}}}
\title{Homework 1}
\author{Satvik Prasad}
\date{}
\newtheorem{lemma}{Lemma}[section]
\fancypagestyle{plain}{%
  \fancyhf{}
  \fancyhead[L]{\theauthor}
  \fancyhead[R]{\thedate}
}
\pagestyle{plain}
\makeatletter
\def\@maketitle{%
  \newpage
  \null
  \vskip 1em%
  \begin{center}%
    \let \footnote \thanks
    {\LARGE \@title \par}%
    \vskip 1em%
  \end{center}%
  \par
\vskip 1em}
\makeatother
\usepackage{cmbright}

\newcommand{\mat}[1]{\mathbf{#1}}

% ---------- Probability theory ----------
% Core operators (optional argument-free; use as \E{X}, \Var{X}, \Cov{X}{Y})
\DeclareMathOperator{\EOp}{\mathbb{E}}
\DeclareMathOperator{\VarOp}{Var}
\DeclareMathOperator{\CovOp}{Cov}
\DeclareMathOperator{\CorrOp}{Corr}

\renewcommand{\P}[1]{\mathbb{P}\left(#1\right)}        % \P{A}
\newcommand{\E}[1]{\EOp\left[#1\right]}                % \E{X}
\newcommand{\Var}[1]{\VarOp\left(#1\right)}            % \Var{X}
\newcommand{\Cov}[2]{\CovOp\left(#1,\,#2\right)}       % \Cov{X}{Y}
\newcommand{\Corr}[2]{\CorrOp\left(#1,\,#2\right)}     % \Corr{X}{Y}
\newcommand{\SD}[1]{\operatorname{SD}\left(#1\right)}  % \SD{X}

% Subscripted versions, e.g. \Esub{\theta}{X}, \Psub{\theta}{A}
\newcommand{\Esub}[2]{\EOp_{#1}\left[#2\right]}
\newcommand{\Psub}[2]{\mathbb{P}_{#1}\left(#2\right)}

% Conditioning, indicators, independence
\newcommand{\given}{\,\middle|\,}                      % \E{X \given Y}
\newcommand{\ind}[1]{\mathbbm{1}\left\{#1\right\}}     % \ind{X > 0}
\newcommand{\indep}{\perp\!\!\!\perp}                  % X \indep Y
\newcommand{\iid}{\overset{\text{iid}}{\sim}}          % X_i \iid \Normal(0,1)

% Convergence
\newcommand{\convp}{\xrightarrow{\;p\;}}               % in probability
\newcommand{\convd}{\xrightarrow{\;d\;}}               % in distribution
\newcommand{\convas}{\xrightarrow{\text{a.s.}}}        % almost surely
\newcommand{\convLp}[1]{\xrightarrow{L^{#1}}}          % in L^p

% Probability spaces
\newcommand{\F}{\mathcal{F}}                           % sigma-algebra
\newcommand{\B}{\mathcal{B}}                           % Borel sigma-algebra
\newcommand{\pspace}{(\Omega, \F, \mathbb{P})}

% Common distributions
\newcommand{\Normal}{\mathcal{N}}
\newcommand{\Bern}{\operatorname{Bernoulli}}
\newcommand{\Bin}{\operatorname{Binomial}}
\newcommand{\Pois}{\operatorname{Poisson}}
\newcommand{\Unif}{\operatorname{Uniform}}
\newcommand{\Exp}{\operatorname{Exp}}
\newcommand{\Geom}{\operatorname{Geometric}}
\newcommand{\Gam}{\operatorname{Gamma}}
\newcommand{\Beta}{\operatorname{Beta}}

% Number sets
\newcommand{\R}{\mathbb{R}}
\newcommand{\N}{\mathbb{N}}
\newcommand{\Z}{\mathbb{Z}}

\begin{document}
\begin{enumerate}
  \item
    \begin{enumerate}
      \item \textbf{Max flow.}
        \begin{enumerate}
          \item The flow on each edge is given below.

            \begin{center}
              \begin{tabular}{c|c|c|c}
                Edge & Flow & Capacity & Full? \\ \hline
                $s \to b$ & 4 & 5 & \\
                $s \to c$ & 4 & 8 & \\
                $b \to e$ & 2 & 2 & yes \\
                $b \to d$ & 2 & 2 & yes \\
                $c \to d$ & 2 & 2 & yes \\
                $c \to f$ & 2 & 3 & \\
                $d \to e$ & 6 & 8 & \\
                $d \to f$ & 0 & 2 & \\
                $f \to h$ & 2 & 2 & yes \\
                $h \to d$ & 2 & 3 & \\
                $e \to g$ & 8 & 10 & \\
                $g \to t$ & 8 & 10 &
              \end{tabular}
            \end{center}

          \item The edges at full capacity are $b \to e$, $b \to d$,
            $c \to d$, and $f \to h$.

          \item The total amount of flow from $s$ to $t$ is $8$.
        \end{enumerate}

      \item \textbf{Min cut.}
        \begin{enumerate}
          \item $W = \{s, b, c, f\}$.

          \item The cut edges are $b \to e$, $b \to d$, $c \to d$,
            and $f \to h$.

          \item The total weight of the cut is $2 + 2 + 2 + 2 = 8$.
        \end{enumerate}
    \end{enumerate}

  \item We are given a directed graph $G = (V, E)$ where every edge
    has capacity 1, and two vertices $s, t.$ We've already solved
    $st$-max flow $f \in \mathbb R^{|E|}.$ We assume that $f$ is
    integral (algorithm from lecture returns an integral flow). Now,
    we insert a new edge $e = (u, v)$ with capacity $1.$ We want to
    find an $\mathcal O(|E|)$ time algorithm to find the new max-flow
    on this graph.

    \begin{enumerate}
      \item
        \begin{algorithmic}[1]
          \Procedure{MaxFlowInsertingEdge}{$G, f, e$}
          \State $G \gets (V(G), E(G) \cup \{e\})$ \Comment{Add $e$ to $G$}
          \State $f(e) \gets 0$ \Comment{Initialize flow of $e$ to 0}
          \vspace{3mm}
          \State $G_f \gets \text{Augment}$($G, f$) \Comment{Generate
          augmented graph $G_f$}
          \State $p \gets \text{DFSPath}$($G_f, s, t$) \Comment{Find
          an $st$-path}

          \vspace{3mm}
          \If{$p \equiv \texttt{nil}$} \Comment{No such st-path
          exists, flow already optimal}
          \State \Return {$f$}
          \EndIf

          \vspace{3mm}

          \For{$e \in p$} \Comment{Push flow along $p$, after we're
          done flow is optimal}
          \If{$e$ \text{forward}}
          \State $f(e) \gets f(e) + 1$
          \Else
          \State $f(e) \gets f(e) - 1$
          \EndIf
          \EndFor

          \State \Return{$f$}
          \EndProcedure
        \end{algorithmic}

      \item \textbf{Correctness} We assume that $f$ is an integral
        flow. Assign $f(e) = 0.$ Then, every $st$-path in the
        augmented graph $G_f$ with edge $e$ added must go through
        $e,$ since otherwise the augmented graph of $G \setminus
        \{e\}$ has an st-path, and is thus not a max-flow.

        Let $W$ be the set of reachable vertices from $s$ on the
        augmentation graph $(G \setminus \{e\})_f.$

        \begin{enumerate}
          \item If there is an $st$-path $p$ on $G_f,$ define flow
            $f^\star$ such that
            $$
            f^\star(e) =
            \begin{cases}
              f(e) + 1& e \in p, \ e \text{ forward}  \\
              f(e) - 1& e \in p, \ e \text{ backward}  \\
              f(e) & \text{otherwise}
            \end{cases}.
            $$

            Then, consider the augmented graph $G_{f^\star}.$ The
            only edges that change from $G_f \to G_{f^\star}$ are
            those along $p,$ and these edges are reversed in
            direction. On $G_f,$ there is no outgoing edge from $W
            \to W^c.$ Since edges are only reversed, on $G_{f^\star}$
            for there to be an outgoing edge from $W \to W^c,$ and
            incoming edge must have been reversed.

            Only edges in $p$ were reversed, for $p$ to take an edge
            from $W^c \to W,$ it must be able to take two unique
            edges from $W \to W^c,$ since $s \in W$ and $t \not \in
            W,$ and $p$ is simple (cannot repeat edges). There can
            only ever be a single edge $e$ from $W \to W^c,$ since
            $f$ was maximal on $(G \setminus \{e\})_f,$ and therefore
            this is a contradiction.

            Thus, on $G_{f^\star}$ there is no outgoing edge from $W
            \to W^c,$ meaning $W$ is a min cut, and $f^\star$ is a max flow.

          \item Otherwise, there is no $st$-path on $G_f.$ Then, the
            set of vertices reachable from $s$ is still a min-cut,
            meaning $f$ is a max flow.
        \end{enumerate}

        In either case, our algorithm correctly returns some
        max-flow, meaning our algorithm is correct.

      \item Modifying the graph, generating the augmentation, and
        updating the flow are all computable in $\mathcal O(|E|)$
        time, meaning our algorithm runs in $\mathcal O(|E|).$
    \end{enumerate}

  \item Want to optimize revenue generated. Letting $x_j$ be the
    volume of material $j,$ we get optimization $$\max c^Tx,$$
    where $$c =
    \begin{pmatrix}
      1000 \\
      1200 \\
      12000
    \end{pmatrix}.$$

    Our constraints satisfy
    $$
    Ax \leq
    \begin{pmatrix}
      100 \\
      60
    \end{pmatrix},
    $$
    where $$A =
    \begin{pmatrix}
      2 & 1 & 3 \\
      1 & 1 & 1
    \end{pmatrix}.$$

  \item
    \begin{enumerate}
      \item Our LP is
        $$
        \min c^T x, \ Ax \geq b,
        $$
        where
        \begin{align*}
          c =
          \begin{pmatrix}
            3 \\
            2
          \end{pmatrix}, \ A =
          \begin{pmatrix}
            1 & 3 \\
            2 & 1 \\
            1 & 0 \\
            0 & 1 \\
          \end{pmatrix}, \ b =
          \begin{pmatrix}
            8 \\
            20 \\
            0 \\
            0
          \end{pmatrix},
        \end{align*}
        where $x = (x_A, x_B)^T,$ since we encode the optimization
        constraint, (in terms of total cost), and the constraints for
        the no. of oranges, bananas, and minimum number of each bundle.

      \item Take $x_A = 10,$ $x_B = 0.$

      \item Take a linear combination of the second and fourth
        constraint, with coefficients $3/2$ and $1/2,$ yielding $(3,
        2) x \geq 30,$ meaning the objective has lower-bound $30.$
    \end{enumerate}

  \item
    \begin{enumerate}
      \item Given $f_{uv}$ for each $(u, v) \in E,$ our cost is given by
        $$
        w^T f,
        $$
        where $f_{(u, v)}= f_{uv},$ and this cost is minimized.

      \item For vertex $v,$ we need
        $$
        \sum_{(u, v) \in E} f_{(u, v)} - \sum_{(v, u) \in E} f_{(v, u)} = d_{v},
        $$
        which is a linear constraint.

      \item We need $f \leq c \in \mathbb R^{|E|},$ where $f_{(u, v)} = f_{uv}.$

      \item Similarly, we need $f \geq 0.$
    \end{enumerate}
\end{enumerate}
\end{document}
